\documentclass[12pt]{article}

\usepackage{graphicx}
\usepackage{color}
\usepackage{amsmath,amssymb}
\usepackage{float}
\usepackage{subcaption}

\begin{document}
\newcommand\reportEIGHTc{0.5175}
\newcommand\reportEIGHTx{0.1593}

\newcommand\reportONEangi{34.4}
\newcommand\reportONEangii{41.6}
\newcommand\reportONEangiii{72.6}
\newcommand\reportONEangix{90.8}
\newcommand\reportONEangx{125.2}

\newcommand\reportTWOfwhm{0.0011}
\newcommand\reportTWOlz{130.023}

\newcommand\reportTHREEfwhm{0.0003}
\newcommand\reportTHREElz{452.025}


    \begin{center}
        \large{\bf{Anisotropic flow analysis with two- and multiparticle correlation techniques}}\\
    \end{center}

    \begin{center}
        {\small\today}
    \end{center}

    %1
    \textbf{Question 1:} {\it Plot the long $2\theta/\omega$ scan of sample B in a logarithmic scale. Which peaks can you identify?} 
    \begin{figure}[H]
        \centering
        \includegraphics[width=0.75\textwidth]{fig_01.png}
        \caption{$2\theta/\omega$ wide range scan of sample B (002)}
        \label{fig_01}
    \end{figure}

    The peaks correspond to Bragg reflections of the $ZnO$ film and the $Al_2O_3$ substrate.
    
    \begin{itemize}
        \item $2\theta = \reportONEangi^{\circ} \rightarrow ZnO$ (002)
        \item $2\theta = \reportONEangii^{\circ} \rightarrow Al_2O_3$ (006)
        \item $2\theta = \reportONEangiii^{\circ} \rightarrow ZnO$ (004)
        \item $2\theta = \reportONEangix^{\circ} \rightarrow Al_2O_3$ (0012)
        \item $2\theta = \reportONEangx^{\circ} \rightarrow ZnO$ (006)
    \end{itemize}

    %2
    \vspace{0.5cm}
    \textbf{Question 2:} {\it Estimate the layer thickness of sample A by using Equation (41).} 
    \begin{figure}[H]
        \centering
        \includegraphics[width=0.75\textwidth]{fig_02.png}
        \caption{$2\theta/\omega$ scan of sample A (002) and Gaussian fit}
        \label{fig_02}
    \end{figure}

    According to equation (41): 
    \begin{equation}
        l_z=\frac{0.9\lambda}{\Delta(2\theta)\cos\theta}
        \label{eq41}
    \end{equation}

    and $FWHM \approx \reportTWOfwhm$ from Gaussian fit, layer thickness $\reportTWOlz$ is estimated as 130 nm.

    %3
    \vspace{0.5cm}
    \textbf{Question 3:} {\it Determine the layer thickness of sample B by using Equations (41) and (42). Compare the obtained values! Conclusion? } 
    \begin{figure}[H]
        \centering
        \includegraphics[width=0.75\textwidth]{fig_03.png}
        \caption{$2\theta/\omega$ scan of sample B (002) and Gaussian fit}
        \label{fig_03}
    \end{figure}

    According to equation (41) (equation (\ref{eq41}) in this report) and $FWHM \approx \reportTHREEfwhm$ from Gaussian fit, layer thickness $\reportTHREElz$ is estimated as 452 nm.
    Equation (42):

    \begin{equation}
        l_z=\frac{\lambda\sin\theta}{\Delta(\theta)\sin 2\theta}
        \label{eq42}
    \end{equation}
    
    requires the distance between secondary maxima. However, in the measured $2\theta/\omega$ scan of sample B, only a single peak is visible and no clear secondary maxima can be identified. Therefore, equation (42) cannot be applied to this dataset.

    This suggests that the thickness estimate from equation (41) is the only accessible one for sample B from the present measurement.

    %4
    \vspace{0.5cm}
    \textbf{Question 4:} {\it Estimate the screw-type dislocation densities of samples A and B and their lateral crystallite size. Compare the structural quality of both samples.} 
    By assuming that the broadening of rocking curve is only caused by limited size, lower limit for the lateral crystallite size $l_\parallel$ can be estimated by:
    \begin{equation}
        l_\parallel = \frac{0.9 \lambda}{\Delta \omega \sin \theta}
        \label{eq43}
    \end{equation}

   \begin{figure}[H]
        \centering
        \begin{subfigure}{0.48\textwidth}
            \centering
            \includegraphics[width=0.75\textwidth]{fig_04_A.png}
            \label{fig_04_A}
        \end{subfigure}
        \begin{subfigure}{0.48\textwidth}
            \centering
            \includegraphics[width=0.75\textwidth]{fig_04_B.png}
            \label{fig_04_B}
        \end{subfigure}
        \caption{$\omega$ scan of sample A and B (002) and Gaussian fit}
        \label{fig_04}
    \end{figure}

    With FWHM 0.003 for sample A and 0.00016 for sample B, lateral sizes $l_\parallel$ are 149 nm and 2854 nm for sample A and B respectively.
    
    By assuming that the broadening of the rocking curves result only from tilt and twist, equation (45) suggests the relation between screw-type dislocation density $\rho_{screw}$ and FWHM $\Delta \omega_{002}$:
    \begin{equation}
        \rho_{screw} = \frac{\Delta \omega_{002}^2}{4.35 |\mathbf{b}_{screw}|^2}
        \label{eq45}
    \end{equation}
    with the Burger vector $\mathbf{b}_{screw}=[001]$ which equals to lattice constance $c = 5.2024 \AA$. $\rho_{screw}$ are $8.38 \times 10^{12} m^{-2}$ and $2.29 \times 10^{10} m^{-2}$. 
    
    Therefore, sample B has a much larger lateral crystallite size and a much lower screw-type dislocation density than sample A, indicating a better structural quality.

    %5
    \vspace{0.5cm}
    \textbf{Question 5:} {\it Estimate the edge-type dislocation density of sample A.} 

    \begin{figure}[H]
        \centering
        \includegraphics[width=0.75\textwidth]{fig_05.png}
        \caption{$\omega$ scan of sample A (101) and Gaussian fit}
        \label{fig_05}
    \end{figure}


    By assuming that the broadening of the rocking curves result only from tilt and twist, equation (46) suggests the relation between edge-type dislocation density $\rho_{edge}$ and FWHM $\Delta \omega_{101}$:
    \begin{equation}
        \rho_{edge} = \frac{\Delta \omega_{101}^2}{4.35 |\mathbf{b}_{edge}|^2}
        \label{eq46}
    \end{equation}
    with the Burger vector $\mathbf{b}_{edge}=[110]$ which equals to lattice constance $a = 3.2475 \AA$. However, no broadening is identified in this dataset, which gives 0 FWHM. This suggests that sample A has zero edge-type dislocation density.

    %6
    \vspace{0.5cm}
    \textbf{Question 6:} {\it Illustrate graphically which lattice planes are associated with the (112) and the (113) reflexes respectively. Verify the epitaxial relation of ZnO on a c-plane sapphire substrate.} 

    \begin{figure}[H]
        \centering
        \includegraphics[width=0.75\textwidth]{fig_06.png}
        \caption{$\phi$ scan of sample A $ZnO$ (112) and $Al_2O_3$ (113)}
        \label{fig_06}
    \end{figure}

    The $\phi$ scans of the ZnO (112) and $Al_2O_3$ (113) reflections show six peaks separated by $60^\circ$, indicating the hexagonal symmetry of both the film and the substrate. The ZnO peaks are shifted by approximately $30^\circ$ with respect to the sapphire peaks.

    %7
    \vspace{0.5cm}
    \textbf{Question 7:} {\it Determine the peak center of the ZnO (205) reflex. Determine the a- and c-lattice constants of sample B. Which mechanisms can you identify for the broadening of its (205) reflex? Calculate the lattice mismatch by assuming that the ZnO film on the sapphire substrate is fully relaxed and by taking into account the formation of a coincidence lattice. } 

    \begin{figure}[H]
        \centering
        \includegraphics[width=0.75\textwidth]{fig_07.png}
        \caption{RSM of sample B $ZnO$ (205)}
        \label{fig_07}
    \end{figure}

    Lattice constants a and c can be calculated by equation 47 and 48:

    \begin{equation}
        a = \frac{1}{q_\parallel (rlu)} \frac{\lambda}{\sqrt{3}} h\\
        \label{eq47}
    \end{equation}

    \begin{equation}
        c = \frac{1}{q_\perp (rlu)} \frac{\lambda}{2} l
        \label{eq48}
    \end{equation}

    where $q_\parallel (rlu)$ and $q_\perp (rlu)$ are scattering vector q components of the peak center in reciprocal lattice units. Experimentally, a = 0.3255 nm, c = 0.3198 nm.

    The broadening results from finite lateral crystallite size and tilt of the crystallites.

    The lattice mismatch is defined as:

    \begin{equation}
        \frac{\Delta a}{a} = \frac{a_{substrate}-a^{relax}_{film}}{a_{substrate}} = \frac{a_{Al_2O_3}-a'_{ZnO}}{a_{Al_2O_3}}
        \label{eq4}
    \end{equation}


    \begin{figure}[H]
        \centering
        \includegraphics[width=0.75\textwidth]{fig_07_fromguide.png}
        \caption{Figure 9 from user guide. Coincidence lattice of $ZnO$ film and $Al_2O_3$ substrate by $30^\circ$ rotation.}
        \label{fig_07_fromguide}
    \end{figure}

    Considering coincidence lattice, with $a'_{ZnO} = 2a_{ZnO}\cos 30 ^\circ = 5.638 \AA$ from experiment, and $a_{Al_2O_3} = 4.7577 \AA$ from appendix B, lattice mismatch is 18.5\%. Relation between two lattice constants a is shown in \ref{fig_07_fromguide}.

    %8
    \vspace{0.5cm}
    \textbf{Question 8:} {\it Determine the c-lattice constant of sample C. Why was the 006-reflex instead of the (002) reflex used for this purpose? ...Use this relation with the above determined value for cZnO to calculate the Mg-content of sample C. Estimate the error of the obtained result by assuming that the Equation (49) is accurate.}

    The (006) reflection was used because higher-order reflections appear at larger diffraction angles $\theta$, which improves the accuracy in determining the lattice constant. From equation 39:

    \begin{equation}
        \frac{\Delta \lambda}{\lambda} = \cot \theta \cdot \theta + \frac{\Delta d_{hkl}}{d_{hkl}}
        \label{eq39}
    \end{equation}

    one can see that higher $\theta$ suppresses $\Delta \theta$ term.

    \begin{figure}[H]
        \centering
        \includegraphics[width=0.75\textwidth]{fig_08.png}
        \caption{$2\theta/\omega$ scan of sample C (006) in unit of $q_{\perp}$}
        \label{fig_08}
    \end{figure}
     
    \ref{eq48}
    \begin{equation}
        c_{Zn_{1-x}Mg_x O} = c_{ZnO} - x \cdot 0.17 \AA
        \label{eq49}
    \end{equation}



    


\end{document}